\documentclass[../DoAn.tex]{subfiles}
\begin{document}

Khi phát triển các hệ thống phần mềm, một quy trình phát triển có vai trò quan trọng trong việc đảm bảo chất lượng, hiệu quả và tiến độ của dự án, đặc biệt là trong các dự án lớn, phức tạp với sự tham gia của một nhóm. Trong lịch sử của ngành công nghiệp phần mềm, đã có nhiều quy trình phát triển phần mềm được đề xuất và áp dụng, với cac ưu điểm và nhược điểm khác nhau.

Chương 1 sẽ trình bày tổng quan về các quy trình phát triển phầm mềm phổ biến như mô hình thác nước (Waterfall), mô hình tăng trường (Iterative), mô hình xoắn ốc (Spiral), mô hình Agile. Đôi với mỗi mô hình, nhóm sẽ trình bày các đặc điểm chính, ưu điểm và nhược điểm của mô hình đó. Đặc biệt, nhóm sẽ tập trung vào mô hình Agile, một trong những quy trình phát triển phần mềm phổ biến nhất hiện nay, với các phương pháp như Scrum, Kanban, XP (Extreme Programming), v.v.

\end{document}